\documentclass[11pt]{article}
\usepackage[margin=1in]{geometry}
\usepackage{amsmath,amssymb,amsthm}
\usepackage{booktabs}
\usepackage{enumitem}
\usepackage[hidelinks]{hyperref}

\newtheorem{proposition}{Proposition}
\newtheorem{corollary}{Corollary}
\theoremstyle{remark}
\newtheorem*{remark}{Remark}

\newcommand{\W}{\mathcal{W}}
\newcommand{\J}{\mathcal{J}}
\newcommand{\Lh}{\widehat{L}}
\newcommand{\Lt}{\widetilde{L}}
\newcommand{\CVaR}{\operatorname{CVaR}}

\title{Per-subproblem Benders cuts and the CVaR link}
\author{Stochastic\_CapExpansion}
\date{November 2026}

\begin{document}
\maketitle

\noindent This note compares the master problem as it is built today, with one optimality cut per
scenario, against the proposed version, with one cut per (scenario, time block) subproblem and the
CVaR auxiliary variables tied to the block cost variables. It shows that the proposed master gives a
valid lower bound, never a looser one, and that for one block per scenario it is the same problem as
today's with fewer rows.

\section{Setting and notation}

\begin{itemize}[leftmargin=*]
  \item $x \in X$: first-stage decisions (resource capacities and line expansions) with investment
        cost $c^\top x$. $X$ is the box and fixed-capacity set the master already imposes.
  \item $\omega \in \W$: operational scenarios $\omega = (s,f,k)$ (demand, gas, weather) with
        probabilities $p_\omega > 0$, $\sum_{\omega} p_\omega = 1$.
  \item $b = 1,\dots,B$: time blocks of a scenario (``Subproblems per scenario''). Each block is a run
        of whole representative periods, and no constraint links two periods.
  \item $\theta_{\omega,b}(x) \ge 0$: optimal operating cost of block $b$ of scenario $\omega$ at
        capacity $x$. It is the value of an LP whose right-hand side is affine in $x$, so it is convex
        and piecewise linear in $x$. Because blocks don't interact, the scenario cost is
        \begin{equation}
          \theta_\omega(x) = \sum_{b=1}^{B} \theta_{\omega,b}(x).
        \end{equation}
  \item $\Omega \in [0,1]$: risk weight (the code's \texttt{risk\_aversion\_weight}), and
        $\Psi \in (0,1]$: tail fraction (\texttt{Value-at-Risk percent}).
  \item $\J$: Benders iterates solved so far. Iterate $j$ solves every block at $x^j$ and returns
        $\theta_{\omega,b}(x^j)$ and a subgradient $\beta^j_{\omega,b}$ (the capacity duals). The block
        tangent is
        \begin{equation}
          \ell^j_{\omega,b}(x) = \theta_{\omega,b}(x^j) + \beta^{j\top}_{\omega,b}\,(x - x^j)
          \;\le\; \theta_{\omega,b}(x) \qquad \forall x,
        \end{equation}
        by convexity. Rows from a subproblem that did not solve to optimality have a zero coefficient
        on the cost variable and are ignored throughout.
\end{itemize}

\noindent The true risk-adjusted objective is
\begin{equation}
  F(x) = c^\top x + \Omega \sum_{\omega} p_\omega\, \theta_\omega(x)
         + (1-\Omega)\, \CVaR_\Psi\!\big(\theta(x)\big),
\end{equation}
with the Rockafellar--Uryasev form
\begin{equation}
  \CVaR_\Psi(v) = \min_{\zeta \in \mathbb{R}} \Big\{ \zeta + \tfrac{1}{\Psi}\sum_\omega p_\omega
  \max(0,\, v_\omega - \zeta) \Big\}.
\end{equation}
$\CVaR_\Psi$ is nondecreasing in each $v_\omega$. The risk-neutral case is $\Omega = 1$ with no
$u$, $\zeta$ variables.

\section{The master problem today}

Today the block results are summed before the cut is built, so each iterate gives one cut per
scenario, $\sum_b \ell^j_{\omega,b}$:
\begin{align}
  \text{(M\textsubscript{agg})}\quad
  \min_{x\in X,\ \alpha,\ u,\ \zeta}\quad
    & c^\top x + \Omega \sum_\omega p_\omega\, \alpha_\omega
      + (1-\Omega)\Big[\zeta + \tfrac{1}{\Psi}\sum_\omega p_\omega u_\omega\Big] \notag\\
  \text{s.t.}\quad
    & \alpha_\omega \ge \textstyle\sum_b \ell^j_{\omega,b}(x)
      && \forall \omega,\ j \in \J \label{eq:agg-opt}\\
    & u_\omega \ge \textstyle\sum_b \ell^j_{\omega,b}(x) - \zeta
      && \forall \omega,\ j \in \J \label{eq:agg-cvar}\\
    & u_\omega \ge -\zeta && \forall \omega \label{eq:agg-tail}\\
    & \alpha_\omega \ge 0,\ u_\omega \ge 0. \notag
\end{align}
In the code, \eqref{eq:agg-opt} are the \texttt{optimality\_cut\_*} rows, \eqref{eq:agg-cvar} the
\texttt{cvar\_tail\_cuts\_*} rows (the same tangent again, minus $\zeta$), and \eqref{eq:agg-tail} the
fixed \texttt{cvar\_tail} rows.
Write the aggregated scenario estimate as
\begin{equation}
  \Lt_\omega(x) = \max\Big(0,\ \max_{j\in\J} \sum_{b} \ell^j_{\omega,b}(x)\Big).
\end{equation}
The row $u_\omega \ge -\zeta$ is the CVaR cut for the constant tangent $0$, which is valid because
costs are nonnegative. So the three $u$ rows together say $u_\omega \ge \Lt_\omega(x) - \zeta$.

\section{The proposed master problem}

Each block gets its own cost variable and its own cuts, and the CVaR variables are tied to the sum of
the block variables by one fixed row per scenario:
\begin{align}
  \text{(M\textsubscript{blk})}\quad
  \min_{x\in X,\ \alpha,\ u,\ \zeta}\quad
    & c^\top x + \Omega \sum_\omega p_\omega \sum_b \alpha_{\omega,b}
      + (1-\Omega)\Big[\zeta + \tfrac{1}{\Psi}\sum_\omega p_\omega u_\omega\Big] \notag\\
  \text{s.t.}\quad
    & \alpha_{\omega,b} \ge \ell^j_{\omega,b}(x)
      && \forall \omega,\ b,\ j \in \J \label{eq:blk-opt}\\
    & u_\omega \ge \textstyle\sum_b \alpha_{\omega,b} - \zeta
      && \forall \omega \label{eq:link}\\
    & \alpha_{\omega,b} \ge 0,\ u_\omega \ge 0. \notag
\end{align}
Here \eqref{eq:blk-opt} are the \texttt{optimality\_cut\_*} rows, now one per block, and
\eqref{eq:link} replaces the fixed \texttt{cvar\_tail} rows. It is built once, with the master, and
no CVaR cut rows are added during the iterations.
The block estimate is
\begin{equation}
  \Lh_\omega(x) = \sum_{b=1}^{B} \max\Big(0,\ \max_{j\in\J} \ell^j_{\omega,b}(x)\Big).
\end{equation}
In code, $\sum_b \alpha_{\omega,b}$ is the expression \texttt{scenario\_alpha}, and
\texttt{expected\_alpha} and \texttt{cvar\_term} keep their current meaning.

\section{Results}

For a fixed $x$, let $V_{\text{agg}}(x)$ and $V_{\text{blk}}(x)$ be the optimal values of the two
masters over the remaining variables $(\alpha, u, \zeta)$.

\begin{proposition}[Reduced objectives]\label{prop:reduced}
For every $x \in X$ (and $\Omega < 1$; for $\Omega = 1$ drop the CVaR term),
\begin{align}
  V_{\text{agg}}(x) &= c^\top x + \Omega \sum_\omega p_\omega \Lt_\omega(x)
                       + (1-\Omega)\,\CVaR_\Psi\big(\Lt(x)\big),\\
  V_{\text{blk}}(x) &= c^\top x + \Omega \sum_\omega p_\omega \Lh_\omega(x)
                       + (1-\Omega)\,\CVaR_\Psi\big(\Lh(x)\big).
\end{align}
\end{proposition}

\begin{proof}
Take (M\textsubscript{blk}) with $x$ fixed. The rows involving $\alpha_{\omega,b}$ are its cuts,
$\alpha_{\omega,b} \ge 0$ and the link \eqref{eq:link}. The objective coefficient of
$\alpha_{\omega,b}$ is $\Omega p_\omega \ge 0$, and in \eqref{eq:link} a smaller $\alpha_{\omega,b}$
only relaxes the bound on $u_\omega$. So, for any feasible point, replacing each $\alpha_{\omega,b}$
by its smallest feasible value $a_{\omega,b} = \max(0, \max_j \ell^j_{\omega,b}(x))$ keeps it
feasible and does not raise the objective. We may therefore set $\alpha_{\omega,b} = a_{\omega,b}$,
so $\sum_b \alpha_{\omega,b} = \Lh_\omega(x)$.

With that, the smallest feasible $u_\omega$ for a given $\zeta$ is
$\max(0, \Lh_\omega(x) - \zeta)$, and $u_\omega$ has the nonnegative coefficient
$(1-\Omega)p_\omega/\Psi$, so it takes that value. What remains is
\[
  c^\top x + \Omega\sum_\omega p_\omega \Lh_\omega(x)
  + (1-\Omega)\min_\zeta\Big\{\zeta + \tfrac1\Psi \sum_\omega p_\omega \max\big(0, \Lh_\omega(x)-\zeta\big)\Big\},
\]
and the minimum over $\zeta$ is $\CVaR_\Psi(\Lh(x))$.

(M\textsubscript{agg}) is the same argument with one block: the smallest feasible $\alpha_\omega$ is
$\Lt_\omega(x)$, and the three families of $u$ rows give
$u_\omega \ge \max\big(0, \max(0, \max_j \sum_b \ell^j_{\omega,b}(x)) - \zeta\big)
= \max(0, \Lt_\omega(x) - \zeta)$.
\end{proof}

\begin{proposition}[Valid and never looser]\label{prop:tight}
For every $x$ and every scenario $\omega$,
\begin{equation}
  \Lt_\omega(x) \;\le\; \Lh_\omega(x) \;\le\; \theta_\omega(x).
\end{equation}
Hence $V_{\text{agg}}(x) \le V_{\text{blk}}(x) \le F(x)$ for every $x$, and the optimal value of
(M\textsubscript{blk}) is a lower bound on $\min_x F(x)$ that is at least as tight as today's.
\end{proposition}

\begin{proof}
\emph{Upper inequality.} Each block term satisfies
$\max(0, \max_j \ell^j_{\omega,b}(x)) \le \theta_{\omega,b}(x)$, because every tangent is below
$\theta_{\omega,b}$ and $\theta_{\omega,b} \ge 0$. Summing over $b$ gives
$\Lh_\omega(x) \le \theta_\omega(x)$.

\emph{Lower inequality.} For any iterate $j$ and each block,
$\ell^j_{\omega,b}(x) \le \max(0, \max_{j'} \ell^{j'}_{\omega,b}(x))$, so summing over $b$,
$\sum_b \ell^j_{\omega,b}(x) \le \Lh_\omega(x)$ for every $j$. Also $\Lh_\omega(x) \ge 0$. Taking the
maximum over $j$ and over $0$ gives $\Lt_\omega(x) \le \Lh_\omega(x)$.

The bounds on $V$ follow from Proposition~\ref{prop:reduced}: both the expectation and
$\CVaR_\Psi$ are nondecreasing in the scenario values, and $\Omega, 1-\Omega \ge 0$.
\end{proof}

\noindent The gain is that each block may take its best tangent from a different iterate. For
example, a summer week can use the tangent from iterate 3 and a winter week the one from iterate 7.
Today one iterate has to supply all of a scenario's blocks at once. The lower inequality is strict
whenever no single iterate is best for every block of a scenario at $x$.

\begin{corollary}[One block per scenario]\label{cor:b1}
If $B = 1$, then $\Lh_\omega = \Lt_\omega$ for every $\omega$. So $V_{\text{blk}} = V_{\text{agg}}$, and
(M\textsubscript{blk}) has the same optimal $x$, lower bound and CVaR term as today's master. It
replaces the per-iteration \texttt{cvar\_tail\_cuts} rows with the single fixed row
\eqref{eq:link} per scenario.
\end{corollary}

\begin{proof}
With $B = 1$, $\Lh_\omega(x) = \max(0, \max_j \ell^j_{\omega,1}(x)) = \Lt_\omega(x)$. Since the two
reduced objectives in Proposition~\ref{prop:reduced} are then the same function of $x$, the two
masters have the same minimizers and optimal value.
\end{proof}

\begin{remark}[Why tie CVaR to the block variables instead of keeping CVaR cuts]
CVaR cuts on $u_\omega$ can only use one iterate's tangent per row:
$u_\omega \ge \sum_b \ell^j_{\omega,b}(x) - \zeta$. That bounds $u_\omega$ by $\Lt_\omega(x) - \zeta$
however many block cuts there are, so the CVaR term would stay at today's looser estimate while the
expectation term gets $\Lh$. With the link \eqref{eq:link}, the CVaR term picks up $\Lh$ too
(Proposition~\ref{prop:reduced}), and the master has fewer rows.
\end{remark}

\begin{remark}[$\Omega = 0$, which is the same as today]
When $\Omega = 0$, $\alpha$ appears in the objective only through $u$. For a scenario outside the
tail, the link \eqref{eq:link} is slack, so $\alpha_{\omega,b}$ can take any value at or above its
cuts with the same objective. The optimal $x$, the lower bound
$c^\top x + \CVaR_\Psi(\Lh(x))$ and the feasible set in $x$ are unaffected, since the argument of
Proposition~\ref{prop:reduced} still gives the optimal value. Only the reported \texttt{Alpha}
values can sit above $\Lh$. Today, at $\Omega = 0$, $\alpha$ has zero objective weight and appears only
in its cuts (and in any budget rows), so it is at least as loose. Budget constraints that contain \texttt{expected\_alpha}
only become harder to meet if $\alpha$ rises, so they do not change which $x$ is feasible either.
\end{remark}

\section{What else stays the same}

\begin{itemize}[leftmargin=*]
  \item \textbf{Lower bound} reported by the Benders loop:
        $\Omega\cdot\texttt{expected\_alpha} + (1-\Omega)\cdot\texttt{cvar\_term} + c^\top x$,
        unchanged in form. By Proposition~\ref{prop:reduced} it equals $V_{\text{blk}}$ at the
        master solution.
  \item \textbf{Upper bound}: computed from the true scenario costs $\theta_\omega(x^j)$, as now.
  \item \textbf{GridFill and the cut export} use cuts without the master, so they get the same
        block form. A scenario's lower bound at a blend $\bar x$ becomes
        $\mathrm{LB}_\omega(\bar x) = \sum_b \max_{k \in (\omega, b)} (a_k + g_k^\top \bar x)$
        instead of $\max_{k \in \omega}(a_k + g_k^\top \bar x)$. It is valid for the same reason as
        the upper inequality of Proposition~\ref{prop:tight}.
\end{itemize}

\section{Size of the master}

Per Benders iteration, with $|\W|$ scenarios and $B$ blocks per scenario:

\begin{center}
\begin{tabular}{@{}lcc@{}}
  \toprule
                                   & Today (M\textsubscript{agg}) & Proposed (M\textsubscript{blk}) \\
  \midrule
  Cost variables                   & $|\W|$                      & $B\,|\W|$ \\
  Optimality cut rows / iteration  & $|\W|$                      & $B\,|\W|$ \\
  CVaR cut rows / iteration        & $|\W|$ (risk on)            & $0$ \\
  Fixed CVaR rows                  & $|\W|$ ($u \ge -\zeta$)     & $|\W|$ (link) \\
  \bottomrule
\end{tabular}
\end{center}

\noindent With risk on, $B = 1$ halves the cut rows added per iteration. With $B = 12$ (one per week
on the 16-zone case) each iteration adds $12|\W|$ rows, six times today's $2|\W|$. The cut pool's
pruning and archiving keep the live set small. Fewer iterations should offset part of the growth,
but the ``Cuts retained'' cap may need raising for large runs.

\end{document}
